\chapter{Weight Gain (Menopausal Weight Change)}

\section*{Definition}
Unintentional increase in body weight, typically with a shift toward abdominal (visceral) fat distribution, occurring during the perimenopausal and postmenopausal transition.

\section*{What Happens in Your Body}
Menopausal weight gain results from several mechanisms: declining oestrogen reduces basal metabolic rate by approximately 50-100 kcal per decade; oestrogen deficiency increases visceral adipose tissue deposition by altering lipoprotein lipase activity; insulin sensitivity decreases, promoting fat storage; loss of lean muscle mass (sarcopenia) further reduces resting energy expenditure. Sleep disruption elevates cortisol and ghrelin while reducing leptin, increasing appetite and cravings. Fat distribution shifts from gynaecoid (hip/thigh) to android (abdominal) pattern, increasing cardiometabolic risk.

\section*{Causes}
\begin{itemize}
  \item Oestrogen decline (reduced metabolic rate, increased visceral fat)
  \item Age-related sarcopenia (loss of muscle mass)
  \item Insulin resistance and glucose intolerance
  \item Sleep disruption (altered ghrelin, leptin, and cortisol)
  \item Thyroid dysfunction (hypothyroidism)
  \item Medication side effects (SSRIs, antipsychotics, beta-blockers, corticosteroids)
  \item Reduced physical activity (joint pain, fatigue)
  \item Emotional eating (mood swings, depression, stress)
  \item Cushing syndrome (cortisol excess, rare)
  \item Genetic predisposition to obesity
\end{itemize}

\section*{When to See a Doctor}
See your doctor if:
\begin{itemize}
  \item Symptoms are severe or getting worse over time
  \item Rapid or excessive weight gain (more than 5 kg in 6 months), weight gain with fatigue, cold intolerance, or new onset of metabolic syndrome features
  \item You have other concerning symptoms such as fever, weight loss, or bleeding
\end{itemize}

\section*{Related Symptoms}
\begin{itemize}
  \item Fatigue
  \item Joint pain
  \item Bloating
  \item Shortness of breath
  \item Mood swings
\end{itemize}

\textbf{Important:} If this symptom persists, your doctor will check for serious underlying causes.

\section*{Self-Care / Home Management}
\begin{itemize}
  \item Adopt a structured calorie deficit (300--500 kcal reduction per day) targeting 0.5--1 kg per week, with emphasis on protein (1.2--1.6 g per kg body weight) to preserve lean muscle mass.
  \item Engage in resistance training (two to three sessions per week) combined with 150 minutes of moderate-intensity aerobic exercise weekly to counteract sarcopenia and metabolic slowing.
  \item Prioritise seven to nine hours of quality sleep per night; address sleep disruption with good sleep hygiene and treat contributing factors (hot flashes, RLS, nocturia).
  \item Keep a food and symptom diary to identify emotional eating patterns, carbohydrate cravings, and triggers linked to mood changes.
\end{itemize}

\section*{How Your Doctor Will Evaluate You}
\begin{itemize}
  \item Anthropometric assessment: weight, height, BMI, waist circumference (88 cm or more in women indicates high cardiometabolic risk), and calculation of weight change trajectory.
  \item Laboratory workup: TSH, free T4, fasting glucose, HbA1c, lipid panel, ALT, creatinine, and haemoglobin to screen for hypothyroidism, insulin resistance, and metabolic syndrome.
  \item Evaluate medication list for contributors: SSRIs (paroxetine), antipsychotics (olanzapine), beta-blockers, corticosteroids, and hormonal contraceptives.
  \item Referral to endocrinology if Cushing syndrome is suspected; screen with overnight dexamethasone suppression test or 24-hour urinary free cortisol.
\end{itemize}

\section*{Treatment Options}
\begin{itemize}
  \item Lifestyle intervention is first-line: dietary counselling, exercise prescription (resistance plus aerobic), behavioural therapy, and sleep optimisation; refer to a registered dietitian and exercise physiologist.
  \item Pharmacotherapy for BMI 30 or higher (or 27 or higher with comorbidity): GLP-1 receptor agonists (semaglutide 2.4 mg weekly or liraglutide 3 mg daily) are most effective.
  \item For menopausal weight gain with vasomotor symptoms: consider menopausal hormone therapy, which may attenuate visceral fat accumulation.
  \item Consider anti-obesity medications or bariatric surgery referral for patients with BMI 35 or higher with comorbidities who fail lifestyle and pharmacotherapy.
\end{itemize}

\section*{References}
\begin{thebibliography}{99}
\bibitem{uptodate-weight-gain} Bray GA, et al. Obesity: evaluation and management. In: UpToDate, Post TW (Ed), UpToDate, Waltham, MA; 2025.
\bibitem{nejm-obesity} Gadde KM, et al. "Obesity: Pathophysiology and Management." \emph{N Engl J Med}. 2023;389(11):1021-1033.
\bibitem{lancet-obesity} Bl\"uher M, et al. "Obesity: a complex disease." \emph{Lancet}. 2024;403(10421):496-512.
\bibitem{jam-weight-gain} Apovian CM, et al. "Pharmacologic management of obesity." \emph{JAMA}. 2024;331(7):598-610.
\bibitem{cochrane-weight-gain} M\"uller MJ, et al. "Lifestyle interventions for weight gain prevention." \emph{Cochrane Database Syst Rev}. 2023;(9):CD013654.
\bibitem{harrison-obesity} Longo DL, et al. \emph{Harrison\'s Principles of Internal Medicine}. 21st ed. McGraw-Hill; 2022. Chapter 397: Obesity.
\end{thebibliography}
